\documentclass{report}
\usepackage{amsmath}
\usepackage{amssymb}
\usepackage{hyperref}
\usepackage[margin=1in]{geometry}
\setlength{\parindent}{0pt}

\begin{document}
\thispagestyle{empty}
\begin{center} \LARGE Assignment 3 \end{center}

\section*{Part 1: Mathematical Derivations}

\subsection*{Q1. Fitting an SVM by hand}
(a)
$$\phi(x_1) = \left[ 1, \sqrt{2}x_1, x_1^2\right]^T
            = \left[ 1, (\sqrt{2})(0), (0)^2\right]^T
            = \begin{bmatrix} 1 \\ 0 \\ 0\end{bmatrix}$$
$$\phi(x_2) = \left[ 1, \sqrt{2}x_2, x_2^2\right]^T
            = \left[ 1, (\sqrt{2})(\sqrt{2}), (\sqrt{2})^2\right]^T
            = \begin{bmatrix} 1 \\ 2 \\ 2\end{bmatrix}$$
(b)

The vector that connects \(\phi(x_1)\) to \(\phi(x_2)\) is obtained by

\begin{equation}
\begin{split}
    \phi(x_1) &= w^\perp + \\
    \phi(x_1) -\phi(x_2) &= w^\perp \\
    w^\perp &= \begin{bmatrix} 0 \\ -2 \\ -2\end{bmatrix}
\end{split}
\end{equation}

thus

\begin{equation}
    w = [0,1,1]^T
\end{equation}

is a valid vector that is parallel to optimal \(w\).

(c)
margin = 

\begin{equation}
    \left\|\phi(x_1)-\phi(x_2)\right\| = \sqrt{8}
\end{equation}

which means 

\begin{equation}
    \frac{1}{\left\|w\right\|} = \frac{\sqrt{8}}{2} \implies \left\|w\right\| = \frac{1}{\sqrt{2}}  
\end{equation}

(d)
note, if \(\left\|w\right\| = \frac{1}{\sqrt{2}}\), then \(w = [0,0.5,0.5]^T\)
to find bias we can use

\begin{equation}
    -1 \left( [0, 0.5, 0.5] \begin{bmatrix} 1 \\ 0 \\ 0 \end{bmatrix} + w_0 \right) = 1
\end{equation}

the first term is 0, so \(w_0=-1\).

we can check using 

\begin{equation}
    +1 \left( [0, 0.5, 0.5] \begin{bmatrix} 1 \\ 2 \\ 2 \end{bmatrix} + w_0 \right) = 1
\end{equation}

where we get 

\begin{equation}
    (0+1+1+w_0) = 1 \implies w_o = -1
\end{equation}


(e)
\begin{equation}
    f(\phi(x)) = \text{sign}([0,0.5,0.5] [1,\sqrt{2}x,x^2] + -1)
\end{equation}

the decision in the original space is a quadratic decision boundary. since it is one d, it is moreso deciding along roots of some polynomial in higher dimension.

(f)

\begin{equation}
    \phi(x_1)^T \phi(x_2) = 1
\end{equation}

\begin{equation}
    \begin{split}
        K(x_i, x_j) &= (x_i \cdot x_j + 1)^2\\
        K(0, \sqrt{2}) &= (0 \cdot \sqrt{2} + 1)^2 = 1
    \end{split}
\end{equation}

and so they do match.


\newpage

\subsection*{Q2. Deriving the SVM dual}
(a)

the lagrangian is written as 
\begin{equation}\label{L}
    L(w,b,\alpha) = \frac{1}{2}\left\|w\right\|^2 + \sum_i \alpha_i(1-y_i(w\cdot x_i + b))
\end{equation}


(b)

\begin{equation}\label{1}
    \begin{split}
        \frac{\partial L(w,b,\alpha)}{\partial w} &= \frac{\partial }{\partial w}\left( \frac{1}{2}\left\|w\right\|^2 + \sum_i \alpha_i(1-y_i(w\cdot x_i + b)) \right)\\
        &= \frac{\partial }{\partial w}\left( \frac{1}{2}\left\|w\right\|^2 + \sum_i \alpha_i(y_i(w\cdot x_i )) \right)\\
        &= w - \sum_i \alpha_iy_ix_{i}\\
    \end{split}
\end{equation}
\begin{equation}
    \begin{split}
        \frac{\partial L}{\partial b}&= \frac{\partial}{\partial b}\left( \frac{1}{2}\left\|w\right\|^2 + \sum_i \alpha_i(1-y_i(w\cdot x_i + b)) \right)\\
        &= \frac{\partial}{\partial b}\left( \sum_i \alpha_i(1-y_i(w\cdot x_i + b)) \right)\\
        &= \frac{\partial}{\partial b} \left( -\sum_i \alpha_iy_ib \right)\\
        &= -\sum_i \alpha_iy_i\\
    \end{split}
\end{equation}

if we set these two equations to zero, we get the values we originally wanted. 

\(w\) is a vector representing the weights of the examples by eq~\ref{1}. for each \(x_i\), it is weighted by exactly \(\alpha_i y_i\).


(c)

now that we know \(w=\sum_i \alpha_iy_ix_i\), substituting back into eq~\ref{L}

\begin{equation}\label{3}
    w\cdot w = \sum_i \alpha_iy_ix_i \cdot \sum_i \alpha_iy_ix_i\\ = \sum_{i,j} \alpha_i\alpha_j y_iy_j(x_i\cdot x_j) = \sum_{i,j}\alpha_i\alpha_jy_iy_j\delta_{ij}
\end{equation}
where \(\delta_{ij}\) is the kronecker delta.

and

\begin{equation}\label{2}
    \sum_i \alpha_i(1-y_i(w\cdot x_i+b)) = \sum_{ij} \alpha_i \alpha_j y_iy_j (x_i\cdot x_j) + \sum_i \alpha_i y_i b = \sum_{i,j} \alpha_i\alpha_jy_iy_j \delta_{ij} + 0
\end{equation}

note that eq~\ref{3}, eq~\ref{2} are actually the same. when we input this back into eq~\ref{L},

\begin{equation}
    L(\alpha) = \sum_i \alpha_i - \frac{1}{2}\sum_{i,j} \alpha_i \alpha_j y_iy_j \delta_{i,j}
\end{equation}

the restriction conditions are just from the restrictions in the lagrange multipliers (\(\alpha_n\geq 0\)) and the restriction in the extremum in step b (\(\sum_n \alpha_ny_n = 0\))



(d)

the entire model is only dependent on a small portion of the training data. most points are ignored when \(\alpha_n\) is zero. in soft margin, \(\alpha_n\) is no longer bounded and thus there can exist non margin support vectors (misclassified or insufficiently classified points)



\newpage

\subsection*{Q3. The kernel trick and valid kernels}
(a)

\begin{equation}
    \begin{split}
        K(x, z) &= (x_1 z_1)^2 + (x_2 z_2)^2 + (1)^2 + 2(x_1 z_1)(x_2 z_2) + 2(x_1 z_1)(1) + 2(x_2 z_2)(1)\\
        &= x_1^2 z_1^2 + x_2^2 z_2^2 + 1 + 2 x_1 x_2 z_1 z_2 + 2 x_1 z_1 + 2 x_2 z_2
    \end{split}
\end{equation}

now we can just go entry by entry and define the \(\phi\) map.

\begin{equation}
    \phi(x) = \begin{pmatrix} x_1^2 \\ x_2^2 \\ \sqrt{2} x_1 x_2 \\ \sqrt{2} x_1 \\ \sqrt{2} x_2 \\ 1 \end{pmatrix}
\end{equation}
to verify

\begin{equation}
    K(x, z) = (x \cdot z + 1)^2 = (5 + 1)^2 = 6^2 = 36
\end{equation}

and

\begin{equation}
    \begin{split}
        \phi(1, 2) &= \begin{pmatrix} 
            1^2 \\ 
            2^2 \\ 
            \sqrt{2}(1)(2) \\ 
            \sqrt{2}(1) \\ 
            \sqrt{2}(2) \\ 
            1 
        \end{pmatrix} = 
            \begin{pmatrix} 
            1 \\
            4 \\
            2\sqrt{2} \\
            \sqrt{2} \\
            2\sqrt{2} \\
            1 
        \end{pmatrix}\\
            \phi(3, 1) &= \begin{pmatrix} 
            3^2 \\
            1^2 \\
            \sqrt{2}(3)(1) \\
            \sqrt{2}(3) \\
            \sqrt{2}(1) \\
            1 
        \end{pmatrix} 
            = 
            \begin{pmatrix} 
            9 \\
            1 \\
            3\sqrt{2} \\
            3\sqrt{2} \\
            \sqrt{2} \\
            1 
        \end{pmatrix}
    \end{split}
\end{equation}

\begin{equation}
    \phi(x) \cdot \phi(z) = (1)(9) + (4)(1) + (2\sqrt{2})(3\sqrt{2}) + (\sqrt{2})(3\sqrt{2}) + (2\sqrt{2})(\sqrt{2}) + (1)(1) = 36
\end{equation}
(b)

\begin{equation}
    G =  \begin{pmatrix} 
        1 & 1 & 1 \\ 
        1 & 4 & 1 \\ 
        1 & 1 & 4 
    \end{pmatrix}
\end{equation}

the characteristic polynomial for \(G\) is 

\begin{equation}
    -\lambda^3 + 9\lambda^2 - 21\lambda + 9
\end{equation}

which gives eigenvalues \(3,3+\sqrt{6}, 3-\sqrt{6}\)
which are all positive since \(\sqrt{6}<3\)

(c)

consider \(\phi_K = \begin{pmatrix}
    \phi_{K_1}\\
    \phi_{K_2}
\end{pmatrix}\)

the product of any \(\phi_K(v_1)\cdot\\phi_K(v_2) =K_1(v_1,v_2) + K_2(v_1,v_2)\). 

algebraically, any bilinear psd operator \(K:V\times V\to V\) is closed under addition. 

\begin{equation}
    \langle v, (K_1 + K_2)v \rangle = \langle v, K_1 v \rangle + \langle v, K_2 v \rangle \geq 0
\end{equation}

to answer why RBF is PSD, we note that 

\begin{equation}
    \|x - z\|^2 = \|x\|^2 + \|z\|^2 - 2x \cdot z
\end{equation}

and thus

\begin{equation}
    K(x, z) = \exp(-\gamma \|x\|^2) \exp(-\gamma \|z\|^2) \exp(2\gamma x \cdot z)
\end{equation}

the product of PSD operators is PSD and thus the final term after being taylor expanded is a sum of scaled product of PSD operators (\(x\cdot z \) is PSD)






\end{document}
